% =============================================================================
% 7. LIMITATIONS AND DISCUSSION
% =============================================================================
\section{Limitations and Discussion}
\vspace{-0.5ex}
\label{sec:limitations}

We close by stating four limitations that scope our claims and flag
directions for follow-up work.
(i) \textit{Python-only corpus and evaluation.} The mid-training corpus
and the in-domain agent benchmarks (SWE-Bench-Verified, SWE-Bench-Lite)
are exclusively Python; cross-language evidence comes only indirectly
through FullStackBench-EN (Section~\ref{sec:exp:gen}), and whether the
recipe transfers to Java, C++, or Rust ecosystems is left to future work.
(ii) \textit{Teacher dependency for CoT.} The default recipe relies on
Gemini-3-Preview to synthesize rationales; the CoT-source ablation
(Section~\ref{sec:exp:ablation}) shows self-generated rationales recover
most of the gain, but a fully open-source replication of our best
configuration requires a comparably strong open teacher or acceptance of
the residual gap.
(iii) \textit{Partial cross-base validation.} Our cross-base evidence
comes from a single non-Qwen2.5-Coder configuration (Qwen3-8B with
SWE-Lego), which simultaneously varies the post-training pipeline; the
result should be read as ``not specific to a single
pretraining/post-training combination'' rather than as a guarantee across
all base-model families.
(iv) \textit{Modularity assumption.} Function-aware FIM presupposes
meaningfully modular code; on monolithic scripts, generated code, or
notebooks, the selection pipeline has fewer eligible targets and degrades
to a smaller usable subset, a regime we do not study systematically here.
